% Claim proof file for row claim_3_20 -- "AcyclicNonCollidersBlockable".
% Auto-generated stub by scaffold/solving_current_row.py at row start. A
% manager agent dedicated to writing the TeX proof fills in the body below.
% The proof must:
%   - Stay close to the lecture notes' proof if one exists (always search
%     the LN first for a `\begin{proof}` block immediately following the
%     claim; copy / lightly edit when present).
%   - Otherwise use the LN paradigm and cite prior claims by their `ref`.
%   - Be self-contained enough that a mathematician can follow it without
%     re-reading the LN.
%   - Have no `sorry`, `omitted`, or "obvious" shortcuts.
%
% Compiles standalone via the `subfiles` package.

\documentclass[main]{subfiles}

\begin{document}
% ref: claim_3_20 (proof, with statement restated for readability)
% title: AcyclicNonCollidersBlockable
\def\rowref{claim\_3\_20}
\phantomsection\label{claim_3_20}
%
% Cross-references: see claim_<ref>_statement_<title>.tex in this folder
% for the corresponding statement.

% --- statement (restated) ---------------------------------------------------
\begin{Rem}[AcyclicNonCollidersBlockable]
If $G$ is acyclic then all non-colliders are blockable.
\end{Rem}

% --- proof ------------------------------------------------------------------
\begin{proof}
Let $v_k$ be a non-collider on $\pi$. By def 3.16, $v_k$ is a blockable non-collider on $\pi$ if and only if it is a non-collider on $\pi$ and is \emph{not} an unblockable non-collider on $\pi$. The non-collider clause already holds by hypothesis, so it suffices to show that $v_k$ is not unblockable on $\pi$.

Suppose for contradiction that $v_k$ were an unblockable non-collider on $\pi$. Then, by def 3.16, $k \notin \lC 0, n\rC$ (so both neighbours $v_{k-1}$ and $v_{k+1}$ exist on $\pi$) and the joint at $v_k$ falls into one of the three patterns
\[\begin{array}{rccl}
  \text{ left chain: } & v_{k-1} \hut v_k \hus v_{k+1} & \text{ with } & v_{k-1} \in \Sc^G(v_k),\\
  \text{ right chain: } & v_{k-1} \suh v_k \tuh v_{k+1} & \text{ with } & v_{k+1} \in \Sc^G(v_k),\\
  \text{ fork: } & v_{k-1} \hut v_k \tuh v_{k+1} & \text{ with } & v_{k-1} \in \Sc^G(v_k) \land v_{k+1} \in \Sc^G(v_k).
\end{array}\]
By def 3.5(7), $\Sc^G(v_k) = \Anc^G(v_k) \cap \Desc^G(v_k) \ins \Anc^G(v_k)$, so any neighbour of $v_k$ on $\pi$ landing in $\Sc^G(v_k)$ in particular admits a directed walk into $v_k$.

In the \emph{left chain} pattern the edge $v_{k-1} \hut v_k$ is the directed edge $v_k \tuh v_{k-1}$ in $G$, i.e.\ a directed walk from $v_k$ to $v_{k-1}$ of length $1$. Combined with the directed walk $v_{k-1} \to \cdots \to v_k$ witnessing $v_{k-1} \in \Anc^G(v_k)$, this concatenates to a directed walk from $v_k$ to itself of length at least $1$, contradicting acyclicity of $G$ (def 3.6).

The \emph{right chain} pattern is symmetric: $v_k \tuh v_{k+1}$ is a directed edge $v_k \to v_{k+1}$, and $v_{k+1} \in \Sc^G(v_k) \ins \Anc^G(v_k)$ supplies a directed walk $v_{k+1} \to \cdots \to v_k$, so concatenation again yields a directed walk from $v_k$ to itself of length at least $1$, contradicting def 3.6.

In the \emph{fork} pattern both arms point outward from $v_k$ and both endpoints lie in $\Sc^G(v_k)$, so either side closes a directed cycle through $v_k$ exactly as in the left or right chain case, contradicting def 3.6.

In all three cases we reach a contradiction. Therefore $v_k$ is not an unblockable non-collider on $\pi$, and combined with the non-collider hypothesis we conclude that $v_k$ is a blockable non-collider on $\pi$.
\end{proof}

\end{document}
